
\documentclass[11pt,a4paper]{article}

\usepackage[utf8]{inputenc}
\usepackage[T1]{fontenc}
\usepackage[spanish,es-nodecimaldot]{babel}
\usepackage{amsmath,amssymb,mathtools}
\usepackage{booktabs,longtable,array,tabularx}
\usepackage{geometry}
\usepackage{xcolor}
\usepackage{hyperref}
\usepackage{enumitem}
\usepackage{siunitx}
\usepackage{microtype}

\geometry{margin=2.1cm}
\hypersetup{
    colorlinks=true,
    linkcolor=blue!55!black,
    urlcolor=blue!55!black
}
\sisetup{output-decimal-marker={,}}

\newcommand{\COtwo}{\mathrm{CO_2}}
\newcommand{\Otwo}{\mathrm{O_2}}
\newcommand{\dd}{\mathrm{d}}
\newcommand{\clip}{\operatorname{clip}}
\newcommand{\e}{\mathrm{e}}
\newcommand{\gLh}{\mathrm{g\,L^{-1}\,h^{-1}}}
\newcommand{\gL}{\mathrm{g\,L^{-1}}}
\newcommand{\mgL}{\mathrm{mg\,L^{-1}}}

\title{\textbf{Modelo dinámico de fermentación y CO\textsubscript{2}\\
para estimación de estados}\\
\large Estructura de trabajo 2026: $S/N/X/E$ + CO\textsubscript{2} disuelto}
\author{}
\date{26 de agosto de 2026}

\begin{document}
\maketitle

\begin{abstract}
Este documento reúne, en una notación única, las ecuaciones actualmente implementadas para el
modelo primario de fermentación y para la capa de producción, acumulación y liberación de
CO\textsubscript{2}. El objetivo es servir como base matemática para integrar la medición online
de CO\textsubscript{2} a un futuro estimador de estados.

La conclusión de modelado adoptada aquí es la siguiente:
\begin{enumerate}[label=(\roman*)]
    \item el CO\textsubscript{2} disuelto, $C_d$, debe conservarse como estado dinámico porque
    introduce la memoria entre producción metabólica y caudal gaseoso;
    \item el O\textsubscript{2} se mantiene en este documento como \emph{estado auxiliar del modelo
    completo actual}, porque la producción efectiva de CO\textsubscript{2} depende de su historia,
    pero no se recomienda todavía comprometerlo como estado obligatorio del EKF;
    \item el vector objetivo $[S,N,X,E,C_d]^\mathsf{T}$ es una reducción deseada, pero el modelo
    implementado distingue glucosa $G$ y fructosa $F$. Por tanto, la reducción
    $S=G+F$ requiere una regla adicional para reconstruir el reparto $G/F$.
\end{enumerate}
\end{abstract}

\tableofcontents
\newpage

\section{Alcance, nomenclatura y niveles del modelo}

La cadena conceptual actual es
\[
(X,N,G,F,E,T,u_N)
\longrightarrow r_E
\longrightarrow q_{\mathrm{bio}}
\longrightarrow q_{\mathrm{prod}}
\longrightarrow C_d
\longrightarrow q_{\mathrm{gas}}
\longrightarrow \widehat y_{\COtwo}.
\]

Las cantidades de CO\textsubscript{2} no deben confundirse:
\begin{align*}
q_{\mathrm{bio}} &:\quad \text{producción biológica base de CO\textsubscript{2}},\\
q_{\mathrm{prod}} &:\quad \text{fuente efectiva que entra al reservorio líquido},\\
C_d &:\quad \text{concentración efectiva de CO\textsubscript{2} disuelto},\\
q_{\mathrm{gas}} &:\quad \text{tasa de liberación desde el líquido hacia el gas},\\
\widehat y_{\COtwo} &:\quad \text{salida del modelo comparable con el caudalímetro}.
\end{align*}

En general,
\[
q_{\mathrm{prod}}\neq q_{\mathrm{gas}}.
\]
Sólo cuando el reservorio disuelto está aproximadamente estacionario,
\[
\frac{\dd C_d}{\dd t}\approx 0,
\]
se obtiene
\[
q_{\mathrm{prod}}\approx q_{\mathrm{gas}}.
\]

\section{Vectores de estado: modelo implementado y modelo objetivo}

\subsection{Modelo primario implementado}

El código primario actualmente integra
\[
\mathbf{x}_{\mathrm{up}}=
\begin{bmatrix}
X & X_d & N & G & F & E & Gly
\end{bmatrix}^{\mathsf T},
\]
donde $X$ es biomasa viable, $X_d$ biomasa muerta, $N$ nitrógeno asimilable,
$G$ glucosa, $F$ fructosa, $E$ etanol y $Gly$ glicerol.

Para el bloque de CO\textsubscript{2}, los estados adicionales de la formulación completa actual son
\[
O_2(t),\qquad C_d(t).
\]
Por tanto, una representación completa de referencia es
\[
\mathbf{x}_{\mathrm{ref}}=
\begin{bmatrix}
X & X_d & N & G & F & E & Gly & O_2 & C_d
\end{bmatrix}^{\mathsf T}.
\]

\subsection{Vector objetivo para el estimador}

El objetivo propuesto es
\[
\boxed{
\mathbf{x}_{\mathrm{EKF},5}=
\begin{bmatrix}
S & N & X & E & C_d
\end{bmatrix}^{\mathsf T}}
\]
con
\[
S=G+F.
\]

Sin embargo, esta reducción no es exactamente cerrada con las ecuaciones actuales porque las
cinéticas de glucosa y fructosa son diferentes y la utilización de fructosa contiene inhibición por
glucosa. La ecuación exacta de $S$ puede escribirse, pero aún requiere conocer $G$ y $F$ por separado.
Esta limitación se discute en la Sección~\ref{sec:s_reduction}.

\subsection{Criterio adoptado para O\textsubscript{2}}

Para reproducir exactamente la capa de CO\textsubscript{2} calibrada desde el inicio de la fermentación,
$O_2$ debe mantenerse como variable dinámica:
\[
\boxed{
\mathbf{x}_{\mathrm{EKF},6}=
\begin{bmatrix}
S & N & X & E & C_d & O_2
\end{bmatrix}^{\mathsf T}}
\]
si se decide estimarlo conjuntamente.

No obstante, el $O_2$ actual es un estado macroscópico latente, no una trayectoria validada por
medición continua. Por ello se propone evaluar dos variantes:
\begin{enumerate}
    \item \textbf{modelo aumentado:} conservar $O_2$ como estado auxiliar;
    \item \textbf{modelo reducido post-anaerobiosis:} iniciar el estimador después del agotamiento
    efectivo de $O_2$, usando aproximadamente $f_{\mathrm{ferm}}\simeq 1$ y
    $q_{\mathrm{resp}}\simeq 0$.
\end{enumerate}

\section{Modelo primario de fermentación}

\subsection{Factores de temperatura}

Sea
\[
T_K=T+273.15,
\]
con $T$ en \si{\degreeCelsius}. Se definen los factores tipo Arrhenius
\begin{align}
a_\mu(T) &=
\exp\left[
E_{ac}\frac{T_K-300}{300RT_K}
\right],\\
a_\beta(T) &=
\exp\left[
E_{afe}\frac{T_K-296.15}{296.15RT_K}
\right],\\
a_{K_N}(T) &=
\exp\left[
E_{aK_N}\frac{T_K-293.15}{293.15RT_K}
\right],\\
a_{K_G}(T) &=
\exp\left[
E_{aK_G}\frac{T_K-293.15}{293.15RT_K}
\right],\\
a_{K_F}(T) &=
\exp\left[
E_{aK_F}\frac{T_K-293.15}{293.15RT_K}
\right],\\
a_{K_{iG}}(T) &=
\exp\left[
E_{aK_{iG}}\frac{T_K-293.15}{293.15RT_K}
\right],\\
a_{K_{iE}}(T) &=
\exp\left[
E_{aK_{iE}}\frac{T_K-293.15}{293.15RT_K}
\right].
\end{align}

Las constantes de semisaturación efectivas son
\begin{align}
K_N(T) &= \frac{\mu_0}{s_N}\,a_{K_N}(T),\\
K_G(T) &= \frac{\beta_{G0}}{s_G}\,a_{K_G}(T),\\
K_F(T) &= \frac{\beta_{F0}}{s_F}\,a_{K_F}(T).
\end{align}

Los coeficientes efectivos de inhibición son
\begin{align}
i_G(T)&=\frac{i_G}{a_{K_{iG}}(T)},\\
i_E(T)&=\frac{i_E}{a_{K_{iE}}(T)}.
\end{align}
En el código se añade un $\varepsilon$ numérico pequeño a los denominadores para evitar divisiones
por cero; no representa un proceso físico adicional.

\subsection{Limitaciones e inhibiciones}

Se definen
\begin{align}
f_N &= \frac{N}{N+K_N(T)},\\
f_G &= \frac{G}{G+K_G(T)},\\
f_F &= \frac{F}{F+K_F(T)},\\
I_E &= \frac{1}{1+i_E(T)E},\\
I_{G\rightarrow F} &= \frac{1}{1+i_G(T)G}.
\end{align}

Los factores efectivos usados por el modelo son
\begin{align}
F_X &= a_\mu(T) f_N,\\
F_G &= a_\beta(T) f_G I_E,\\
F_F &= a_\beta(T) f_F I_{G\rightarrow F} I_E.
\end{align}

Así,
\begin{align}
\mu &= \mu_0 F_X,\\
\beta_G &= \beta_{G0}F_G,\\
\beta_F &= \beta_{F0}F_F.
\end{align}

\subsection{Mantenimiento y muerte celular}

La tasa de mantenimiento es
\begin{equation}
m(T)=
m_0
\exp\left[
E_{am}\frac{T_K-293.3}{293.3RT_K}
\right].
\end{equation}

Sea
\[
S_{\mathrm{tot}}=G+F.
\]
La disponibilidad para mantenimiento es
\begin{equation}
f_m=
\frac{S_{\mathrm{tot}}}
{S_{\mathrm{tot}}+S_{m,1/2}},
\qquad
S_{m,1/2}=1.0.
\end{equation}

El umbral térmico de muerte dependiente de etanol es
\begin{equation}
T_d(E)=
-10^{-4}E^3+0.0049E^2-0.1279E+315.89
\quad [\mathrm{K}].
\end{equation}

La tasa de muerte celular se implementa como
\begin{equation}
k_d=
\begin{cases}
K_{d0}
\exp\left[
C_{de}E+
E_{td}\dfrac{T_K-305.65}{305.65RT_K}
\right],
& T_K\geq T_d(E),\\[8pt]
0,
& T_K<T_d(E).
\end{cases}
\end{equation}

\subsection{Ecuaciones diferenciales primarias}

Las ecuaciones implementadas son
\begin{align}
\frac{\dd X}{\dd t}
&=(\mu-k_d)X,
\label{eq:x}\\
\frac{\dd X_d}{\dd t}
&=k_dX,
\label{eq:xd}\\
\frac{\dd N}{\dd t}
&=-q_NF_XX,
\label{eq:n}\\
\frac{\dd G}{\dd t}
&=
-\left[
q_{XG}F_X
+q_{EG}F_G
+m(T)f_m\frac{G}{G+F}
\right]X,
\label{eq:g}\\
\frac{\dd F}{\dd t}
&=
-\left[
q_{XF}F_X
+q_{EF}F_F
+m(T)f_m\frac{F}{G+F}
\right]X,
\label{eq:f}\\
\frac{\dd E}{\dd t}
&=
(\beta_G+\beta_F)X,
\label{eq:e}\\
\frac{\dd Gly}{\dd t}
&=
\left[
\gamma_{G0}F_G+\gamma_{F0}F_F
\right]X.
\label{eq:gly}
\end{align}

La tasa de producción de etanol utilizada por la capa de CO\textsubscript{2} es
\begin{equation}
\boxed{
r_E(t)=\frac{\dd E}{\dd t}
=
\left[\beta_G(t)+\beta_F(t)\right]X(t)}
\qquad [\gLh].
\label{eq:rE}
\end{equation}

\section{Reducción exacta a azúcar total $S=G+F$}
\label{sec:s_reduction}

Definiendo
\[
S=G+F,
\]
se obtiene exactamente
\begin{align}
\frac{\dd S}{\dd t}
&=
\frac{\dd G}{\dd t}
+
\frac{\dd F}{\dd t}\\
&=
-\Big[
(q_{XG}+q_{XF})F_X
+q_{EG}F_G
+q_{EF}F_F
+m(T)f_m
\Big]X.
\label{eq:S_total}
\end{align}

La ecuación~\eqref{eq:S_total} \textbf{no queda cerrada únicamente con $S$}, porque
$F_G$ depende de $G$ y $F_F$ depende simultáneamente de $F$ y de la inhibición por $G$.

Una forma de escribir explícitamente la reducción es introducir
\[
\rho_G(t)=\frac{G(t)}{S(t)},
\qquad
G=\rho_GS,
\qquad
F=(1-\rho_G)S.
\]
Entonces la ecuación de $S$ puede evaluarse si $\rho_G(t)$ se proporciona como una trayectoria
conocida/nominal. \textbf{Esta reconstrucción de $\rho_G$ todavía no forma parte del modelo
reducido validado}; es una decisión pendiente para construir el EKF de cinco estados.

Por esta razón, la implementación más fiel antes de cerrar esa reducción sería trabajar con
\[
[G,F,N,X,E,C_d]
\]
(o añadiendo $O_2$), y sólo después verificar si la agregación $S=G+F$ conserva suficiente
información para el observador.

\section{Pulsos de nitrógeno}

\subsection{Pulso en el modelo primario}

El simulador primario integra las ODE entre eventos y aplica el pulso de N como un salto exacto:
\begin{equation}
\boxed{
N(t_N^+)=N(t_N^-)+\Delta N.
}
\label{eq:N_jump}
\end{equation}

No se ``difumina'' el salto dentro de la ODE primaria.

\subsection{Disponibilidad gradual usada por la capa de CO\textsubscript{2}}

Para evitar que el pulso se traduzca instantáneamente en una respuesta de CO\textsubscript{2}, la
calibración 2026 utiliza una rampa causal
\begin{equation}
f_N(t)=
\clip\left(
\frac{t-t_N}{t_{\mathrm{rise}}},
0,1
\right).
\label{eq:pulse_fraction}
\end{equation}

Sean $(X_0,q_{\mathrm{bio},0})$ las trayectorias contrafactuales sin pulso y
$(X_{+N},q_{\mathrm{bio},+N})$ las trayectorias con pulso. Entonces
\begin{equation}
X_{\mathrm{eff}}
=
X_0+
f_N
\left(
X_{+N}-X_0
\right),
\label{eq:Xeff}
\end{equation}
y
\begin{equation}
q_{\mathrm{bio},N}
=
\left[
1+(g_N-1)f_N
\right]q_{\mathrm{bio},0}
+
f_N
\left(
q_{\mathrm{bio},+N}-q_{\mathrm{bio},0}
\right).
\label{eq:qbioN}
\end{equation}

Esta formulación es la implementación actualmente calibrada. Para un futuro EKF completamente
integrado deberá decidirse si se conserva esta construcción contrafactual o si se reemplaza por una
dinámica explícita de disponibilidad del pulso.

\section{Producción biológica base de CO\textsubscript{2}}

La producción biológica base se obtiene estequiométricamente desde la tasa de etanol:
\begin{equation}
\boxed{
q_{\mathrm{bio}}(t)
=
\alpha_{\COtwo/E}
\max\left[r_E(t),0\right]
}
\label{eq:qbio}
\end{equation}
con
\begin{equation}
\boxed{
\alpha_{\COtwo/E}
=
\frac{44.01}{2(46.07)}
\simeq 0.4776
}
\qquad
\left[
\frac{\mathrm{g\ CO_2}}{\mathrm{g\ etanol}}
\right].
\label{eq:alphaCO2}
\end{equation}

\section{Submodelo de O\textsubscript{2}: estado auxiliar de la formulación actual}

\subsection{Saturación de referencia e condición inicial}

La correlación de trabajo para O\textsubscript{2} a saturación es
\begin{equation}
O_2^\ast(T,E,G,F)
=
8.6
\exp[-0.024(T-20)]
\exp(-0.0020E)
\exp[-0.0010(G+F)]
\qquad [\mgL].
\label{eq:o2sat}
\end{equation}

La condición inicial utilizada en Laboratorio 2026 es
\begin{equation}
\boxed{
O_2(0)=
s_{O2,0}
O_2^\ast(T_0,E_0,G_0,F_0).
}
\label{eq:o2initial}
\end{equation}

\subsection{Consumo de O\textsubscript{2}}

La tasa de consumo es
\begin{equation}
u_{O2}
=
q_{O2,\max}
X_{\mathrm{eff}}
\frac{O_2}{K_{O2}+O_2},
\label{eq:o2uptake}
\end{equation}
con unidades de \si{mg.L^{-1}.h^{-1}}.

En la formulación 2026 no hay reaireación durante el proceso:
\begin{equation}
\boxed{
\frac{\dd O_2}{\dd t}
=
-u_{O2}.
}
\label{eq:o2ode}
\end{equation}

\subsection{Fracción anaerobia y fracción fermentativa}

La fracción anaerobia efectiva es
\begin{equation}
\phi_{\mathrm{ana}}(O_2)
=
\frac{K_{\mathrm{ana}}^h}
{K_{\mathrm{ana}}^h+O_2^h}.
\label{eq:ana}
\end{equation}

La fracción fermentativa es
\begin{equation}
\boxed{
f_{\mathrm{ferm}}
=
f_C+(1-f_C)\phi_{\mathrm{ana}}.
}
\label{eq:ferm_fraction}
\end{equation}

El término respiratorio de CO\textsubscript{2} se calcula como
\begin{equation}
q_{\mathrm{resp}}
=
\frac{44.01}{32.00}
\frac{u_{O2}}{1000}
\qquad [\gLh].
\label{eq:qresp}
\end{equation}

\section{Activación química gradual}

La evidencia química define un intervalo $[t_L,t_U]$:
\begin{itemize}
    \item $t_U$: primera muestra con caída de $G+F$ de al menos \SI{5}{g.L^{-1}}
    o aumento de $E$ de al menos \SI{2}{g.L^{-1}} respecto de la primera muestra;
    \item $t_L$: muestra química inmediatamente anterior.
\end{itemize}

La posición del inicio dentro del bracket es
\begin{equation}
t_s=
t_L+
f_s(t_U-t_L).
\label{eq:ts}
\end{equation}

La duración es
\begin{equation}
\Delta t_{\mathrm{act}}
=
\max\left[
f_d(t_U-t_L),\,
0.25\ \mathrm{h}
\right].
\label{eq:dtact}
\end{equation}

Se define
\begin{equation}
z(t)=
\clip\left(
\frac{t-t_s}{\Delta t_{\mathrm{act}}},
0,1
\right),
\end{equation}
y la activación \emph{smoothstep}
\begin{equation}
\boxed{
a_{\mathrm{chem}}(t)
=
3z^2-2z^3.
}
\label{eq:achem}
\end{equation}

Si el bracket colapsa ($t_U\le t_L$), la implementación usa un escalón en $t_U$.

\section{Fuente total de CO\textsubscript{2}}

La fuente que alimenta el reservorio disuelto es
\begin{equation}
\boxed{
q_{\mathrm{prod}}
=
a_{\mathrm{chem}}(t)
\left[
q_{\mathrm{bio},N}
f_{\mathrm{ferm}}
+
q_{\mathrm{resp}}
\right].
}
\label{eq:qprod}
\end{equation}

Esta ecuación deja explícito que el O\textsubscript{2} modifica la \emph{fuente metabólica}; no
modifica directamente la ley de solubilidad de CO\textsubscript{2}.

\section{Solubilidad efectiva de CO\textsubscript{2}}

La correlación de trabajo es
\begin{equation}
\boxed{
C_{\COtwo}^{\ast}
=
s_{\mathrm{sat}}\,
1.69
\exp[-0.032(T-20)]
\exp(0.0016E)
\exp[-0.0012(G+F)]
}
\qquad [\gL].
\label{eq:csat}
\end{equation}

El parámetro $s_{\mathrm{sat}}$ es una escala calibrada; la correlación no corresponde a una
medición directa de equilibrio en cada fermentación.

\section{CO\textsubscript{2} disuelto como estado dinámico}

La condición inicial usada actualmente es
\begin{equation}
\boxed{
C_d(0)=0.
}
\label{eq:cd0}
\end{equation}

El balance de masa es
\begin{equation}
\boxed{
\frac{\dd C_d}{\dd t}
=
q_{\mathrm{prod}}-q_{\mathrm{gas}}.
}
\label{eq:cdode}
\end{equation}

Esta es la razón principal para tratar $C_d$ como estado: la salida gaseosa depende de la historia
de acumulación y no sólo del estado biológico instantáneo.

\section{Liberación continua líquido--gas}

La fracción de liberación es
\begin{equation}
\psi(C_d,C_{\COtwo}^{\ast})
=
f_0
+
(1-f_0)
\frac{C_d}{C_d+C_{\COtwo}^{\ast}},
\qquad
f_0=0.05.
\label{eq:psi}
\end{equation}

La liberación potencial es
\begin{equation}
\boxed{
q_{\mathrm{gas,pot}}
=
k_{\mathrm{release}}
C_d
\psi(C_d,C_{\COtwo}^{\ast}).
}
\label{eq:qgaspot}
\end{equation}

En tiempo continuo, la representación física de trabajo es
\[
q_{\mathrm{gas}}\approx q_{\mathrm{gas,pot}},
\]
pero la implementación numérica impone una guarda de conservación de masa. Para un paso
$\Delta t$:
\begin{equation}
q_{\mathrm{gas},k}
=
\min\left(
q_{\mathrm{gas,pot},k},
\frac{C_{d,k}}{\Delta t}
+
q_{\mathrm{prod},k-1}
\right),
\label{eq:qgasguard}
\end{equation}
\begin{equation}
C_{d,k+1}
=
\max\left[
C_{d,k}
+
\Delta t
(q_{\mathrm{prod},k-1}-q_{\mathrm{gas},k}),
0
\right].
\label{eq:cd_discrete}
\end{equation}

La función $\min$ no constituye una cinética adicional: sólo evita liberar más masa que la
disponible.

\section{Ecuación de observación: conexión con el caudalímetro}

La salida predicha se expresa como
\begin{equation}
\boxed{
\widehat y_{\COtwo}(t)
=
g_m\,q_{\mathrm{gas}}(t).
}
\label{eq:measurement}
\end{equation}

Por tanto, el caudalímetro no observa directamente $q_{\mathrm{bio}}$ ni
$q_{\mathrm{prod}}$, sino una señal asociada a la liberación gaseosa.

La señal adquirida en SCCM se convierte a tasa específica mediante
\begin{equation}
\boxed{
y_{\gLh}
=
y_{\mathrm{sccm}}
\frac{60}{1000}
\frac{44.01}{22.414}
\frac{1}{V_L}.
}
\label{eq:sccm}
\end{equation}

Aquí $V_L$ es el volumen de líquido en litros. En el preprocesamiento 2026 se corrige además un
offset operacional por corrida antes de la comparación con el modelo.

\section{Sistema dinámico compacto propuesto}

\subsection{Variante completa con O\textsubscript{2}}

Si se conserva la estructura actual desde el inicio de la fermentación, una representación
compacta para el bloque relevante del observador es
\begin{equation}
\boxed{
\begin{aligned}
\dot X &= f_X(X,N,G,F,E,T),\\
\dot N &= f_N(X,N,G,F,E,T,u_N),\\
\dot G &= f_G(X,N,G,F,E,T),\\
\dot F &= f_F(X,N,G,F,E,T),\\
\dot E &= f_E(X,N,G,F,E,T),\\
\dot O_2 &= -u_{O2}(X,O_2),\\
\dot C_d &= q_{\mathrm{prod}}(X,N,G,F,E,O_2,T,u_N)-q_{\mathrm{gas}},\\
y &= g_m q_{\mathrm{gas}}(C_d,G,F,E,T).
\end{aligned}
}
\label{eq:full_compact}
\end{equation}

\subsection{Variante reducida objetivo $[S,N,X,E,C_d]$}

El objetivo final deseado es
\begin{equation}
\boxed{
\mathbf{x}=
[S,N,X,E,C_d]^\mathsf{T}
}
\end{equation}
con
\begin{equation}
\boxed{
\begin{aligned}
\dot S &= f_S(S,N,X,E,T;\rho_G),\\
\dot N &= f_N(S,N,X,E,T,u_N;\rho_G),\\
\dot X &= f_X(S,N,X,E,T;\rho_G),\\
\dot E &= f_E(S,N,X,E,T;\rho_G),\\
\dot C_d &= q_{\mathrm{prod}}-q_{\mathrm{gas}},\\
y &= g_m q_{\mathrm{gas}}.
\end{aligned}
}
\label{eq:reduced_compact}
\end{equation}

El símbolo $\rho_G$ recuerda que aún debe definirse cómo recuperar la partición glucosa/fructosa
a partir de $S$. No debe ocultarse esta dependencia al formalizar el EKF.

Si además se elimina $O_2$, la simplificación sólo es consistente después de justificar
\[
f_{\mathrm{ferm}}\simeq 1,
\qquad
q_{\mathrm{resp}}\simeq 0,
\]
o de reemplazar el submodelo de O\textsubscript{2} por una función conocida.

\section{Parámetros del modelo primario}

\subsection{Parámetros cinéticos}

\begin{longtable}{p{2.6cm}p{8.4cm}p{3.0cm}}
\toprule
\textbf{Parámetro} & \textbf{Interpretación en la implementación} & \textbf{Estatus}\\
\midrule
\endfirsthead
\toprule
\textbf{Parámetro} & \textbf{Interpretación en la implementación} & \textbf{Estatus}\\
\midrule
\endhead
$\mu_0$ & Coeficiente base de crecimiento específico & calibrado upstream\\
$s_N$ & Parametrización de $K_N=\mu_0/s_N$ & fijado en los drivers actuales\\
$q_N$ & Coeficiente de consumo de nitrógeno asociado al crecimiento & calibrado upstream\\
$q_{XG}$ & Consumo de glucosa asociado al crecimiento & fijado en los drivers actuales\\
$q_{XF}$ & Consumo de fructosa asociado al crecimiento & fijado en los drivers actuales\\
$\beta_{G0}$ & Coeficiente base de producción de etanol desde glucosa & calibrado upstream\\
$s_G$ & Parametrización de $K_G=\beta_{G0}/s_G$ & fijado en los drivers actuales\\
$\beta_{F0}$ & Coeficiente base de producción de etanol desde fructosa & calibrado upstream\\
$s_F$ & Parametrización de $K_F=\beta_{F0}/s_F$ & fijado en los drivers actuales\\
$q_{EG}$ & Consumo de glucosa asociado a la vía de producción de etanol & calibrado upstream\\
$q_{EF}$ & Consumo de fructosa asociado a la vía de producción de etanol & calibrado upstream\\
$i_G$ & Coeficiente base de inhibición de utilización de fructosa por glucosa & calibrado upstream\\
$i_E$ & Coeficiente base de inhibición por etanol & calibrado upstream\\
$K_{d0}$ & Prefactor de muerte celular & calibrado upstream\\
$m_0$ & Coeficiente base de mantenimiento & fijado en los drivers actuales\\
$\gamma_{G0}$ & Coeficiente de formación de glicerol asociado a glucosa & calibrado upstream\\
$\gamma_{F0}$ & Coeficiente de formación de glicerol asociado a fructosa & calibrado upstream\\
\bottomrule
\end{longtable}

\subsection{Valores upstream utilizados por los drivers de Laboratorio 2026}

Los valores siguientes son los que alimentan la capa de CO\textsubscript{2} actual. El cargador
parte de los valores por defecto y los sobrescribe con el ajuste específico de cada matriz.

\begin{longtable}{lrrl}
\toprule
\textbf{Parámetro} & \textbf{MN} & \textbf{MS} & \textbf{Origen}\\
\midrule
\endfirsthead
\toprule
\textbf{Parámetro} & \textbf{MN} & \textbf{MS} & \textbf{Origen}\\
\midrule
\endhead
$\mu_0$ & 0.0787279 & 0.133609 & ajuste upstream\\
$s_N$ & 18.0 & 18.0 & valor base mantenido\\
$q_N$ & 0.0210804 & 0.0384848 & ajuste upstream\\
$q_{XG}$ & 0.1125 & 0.1125 & valor base mantenido\\
$q_{XF}$ & 0.1125 & 0.1125 & valor base mantenido\\
$\beta_{G0}$ & 0.428878 & 1.329860 & ajuste upstream\\
$s_G$ & 0.0300 & 0.0300 & valor base mantenido\\
$\beta_{F0}$ & 3.440399 & 1.820638 & ajuste upstream\\
$s_F$ & 0.0300 & 0.0300 & valor base mantenido\\
$q_{EG}$ & 1.712807 & 2.577082 & ajuste upstream\\
$q_{EF}$ & 5.023017 & 4.147139 & ajuste upstream\\
$i_G$ & 0.0560374 & 0.0263463 & ajuste upstream\\
$i_E$ & 0.0401679 & 0.0554701 & ajuste upstream\\
$K_{d0}$ & 0.00015025 & 0.000561724 & ajuste upstream\\
$m_0$ & 0.0100 & 0.0100 & valor base mantenido\\
$\gamma_{G0}$ & 0.115559 & 0.149529 & ajuste upstream\\
$\gamma_{F0}$ & 0.000100002 & 0.0335320 & ajuste upstream\\
\bottomrule
\end{longtable}

Para MS se utiliza el caso \texttt{historical\_synthetic\_plus\_lot2}. Para MN se utiliza
\texttt{estimability\_historical\_natural/theta.csv}.

\subsection{Constantes fijas de temperatura y muerte}

\begin{longtable}{lrl}
\toprule
\textbf{Constante} & \textbf{Valor} & \textbf{Uso}\\
\midrule
\endfirsthead
\toprule
\textbf{Constante} & \textbf{Valor} & \textbf{Uso}\\
\midrule
\endhead
$R$ & 8.314 & constante de gases en factores Arrhenius\\
$E_{ac}$ & 59453 & factor térmico de crecimiento\\
$E_{afe}$ & 11000 & factor térmico de fermentación\\
$E_{aK_N}$ & 46055 & temperatura de $K_N$\\
$E_{aK_G}$ & 46055 & temperatura de $K_G$\\
$E_{aK_F}$ & 46055 & temperatura de $K_F$\\
$E_{aK_{iG}}$ & 46055 & temperatura de inhibición por glucosa\\
$E_{aK_{iE}}$ & 46055 & temperatura de inhibición por etanol\\
$E_{am}$ & 37681 & temperatura del mantenimiento\\
$C_{de}$ & 0.0415 & dependencia de muerte con etanol\\
$E_{td}$ & 130000 & dependencia térmica de muerte\\
$S_{m,1/2}$ & 1.0 & escala de disponibilidad de azúcar para mantenimiento\\
\bottomrule
\end{longtable}

\section{Parámetros de la capa CO\textsubscript{2}/O\textsubscript{2}}

\subsection{Definiciones y cotas}

\begin{longtable}{p{3.7cm}p{6.5cm}p{2.2cm}p{2.5cm}}
\toprule
\textbf{Parámetro} & \textbf{Significado} & \textbf{Unidad} & \textbf{Cota}\\
\midrule
\endfirsthead
\toprule
\textbf{Parámetro} & \textbf{Significado} & \textbf{Unidad} & \textbf{Cota}\\
\midrule
\endhead
$k_{\mathrm{release}}$ & velocidad efectiva de vaciado de $C_d$ & \si{h^{-1}} & 0.03--25\\
$s_{\mathrm{sat}}$ & escala de la correlación $C_{\COtwo}^{\ast}$ & -- & 0.35--2.50\\
$q_{O2,\max}$ & consumo máximo efectivo de O\textsubscript{2} por biomasa & \si{mg.g^{-1}.h^{-1}} & 0.15--6\\
$s_{O2,0}$ & escala del O\textsubscript{2} inicial respecto a $O_2^\ast$ & -- & 0.05--1.50\\
$t_{\mathrm{rise}}$ & tiempo de disponibilidad gradual del pulso N & \si{h} & 1--72\\
$g_N$ & ganancia de actividad asociada al pulso N & -- & 0.25--4\\
$f_s$ & posición del inicio dentro del bracket químico & -- & 0.01--0.95\\
$f_d$ & duración relativa al ancho del bracket químico & -- & 0.35--2\\
$g_m$ & ganancia entre $q_{\mathrm{gas}}$ y sensor convertido & -- & 0.10--20\\
\bottomrule
\end{longtable}

\subsection{Valores calibrados en Laboratorio 2026}

Un asterisco indica que el valor quedó en una cota activa.

\begin{longtable}{lrr}
\toprule
\textbf{Parámetro} & \textbf{MS} & \textbf{MN}\\
\midrule
\endfirsthead
\toprule
\textbf{Parámetro} & \textbf{MS} & \textbf{MN}\\
\midrule
\endhead
$k_{\mathrm{release}}$ & 15.9857 & 0.424588\\
$s_{\mathrm{sat}}$ & 0.350987$^\ast$ & 0.351389$^\ast$\\
$q_{O2,\max}$ & 0.150000$^\ast$ & 0.150000$^\ast$\\
$s_{O2,0}$ & 0.390441 & 0.195400\\
$t_{\mathrm{rise}}$ & 36.9954 & 64.6587\\
$g_N$ & 1.24358 & 0.250000$^\ast$\\
$f_s$ & 0.010000$^\ast$ & 0.735057\\
$f_d$ & 0.350000$^\ast$ & 1.06391\\
$g_m$ & 2.13427 & 3.11576\\
\bottomrule
\end{longtable}

\subsection{Constantes fijas de la capa O\textsubscript{2}/CO\textsubscript{2}}

\begin{longtable}{lrl}
\toprule
\textbf{Constante} & \textbf{Valor} & \textbf{Significado}\\
\midrule
\endfirsthead
\toprule
\textbf{Constante} & \textbf{Valor} & \textbf{Significado}\\
\midrule
\endhead
$K_{O2}$ & \SI{0.25}{mg.L^{-1}} & semisaturación del consumo de O\textsubscript{2}\\
$K_{\mathrm{ana}}$ & \SI{0.75}{mg.L^{-1}} & O\textsubscript{2} donde $\phi_{\mathrm{ana}}=0.5$\\
$h$ & 2 & exponente Hill de la transición anaerobia\\
$f_C$ & 0.08 & piso fermentativo tipo Crabtree\\
$f_0$ & 0.05 & piso de liberación subsaturada de CO\textsubscript{2}\\
$M_{\COtwo}$ & \SI{44.01}{g.mol^{-1}} & masa molar de CO\textsubscript{2}\\
$M_E$ & \SI{46.07}{g.mol^{-1}} & masa molar de etanol\\
$M_{O2}$ & \SI{32.00}{g.mol^{-1}} & masa molar de O\textsubscript{2}\\
$V_m^\circ$ & \SI{22.414}{L.mol^{-1}} & volumen molar estándar usado en conversión SCCM\\
\bottomrule
\end{longtable}

\section{Criterio de calibración de la capa CO\textsubscript{2}}

Esta sección no forma parte de las ODE, pero documenta cómo se obtuvieron los parámetros anteriores.

Para el experimento $i$ y punto $j$:
\begin{equation}
r_{ij}
=
\frac{
g_mq_{\mathrm{gas}}(t_{ij})-y_{ij}
}{
\sigma_i\sqrt{n_i}
},
\qquad
\sigma_i=
\max(0.04,\,0.10y_{i,\max}).
\label{eq:residual}
\end{equation}

Los puntos bajo el límite operacional se tratan como censurados a izquierda:
\begin{equation}
r_{ij}^{\mathrm{cens}}
=
\frac{
\max(\widehat y_{ij}-LOD_{ij},0)
}{
\sigma_i\sqrt{n_i}
}.
\label{eq:censored}
\end{equation}

Además se incorporan penalizaciones por error de onset y por tiempo del máximo posterior al pulso.
El ajuste utiliza mínimos cuadrados robustos con pérdida \texttt{soft\_l1}.

\section{Qué se recomienda llevar al EKF}

\subsection{CO\textsubscript{2} disuelto}

Se recomienda \textbf{sí incluir $C_d$ como estado}. El argumento es estructural:
\[
\dot C_d=q_{\mathrm{prod}}-q_{\mathrm{gas}}
\]
y la medición es
\[
y=g_mq_{\mathrm{gas}}.
\]
Eliminar $C_d$ equivale a aproximar la medición gaseosa por la producción metabólica instantánea,
supuesto que fue precisamente el que produjo problemas de onset y memoria en las comparaciones
anteriores.

\subsection{O\textsubscript{2}}

Se recomienda \textbf{mantener $O_2$ en el modelo matemático de referencia, pero no decidir aún
que deba ser un estado estimado por el EKF}. Las razones son:
\begin{enumerate}
    \item la formulación actual de $q_{\mathrm{prod}}$ sí depende de $O_2$;
    \item no existe una trayectoria continua de O\textsubscript{2} usada para validar directamente ese estado;
    \item varios parámetros de la capa O\textsubscript{2}/CO\textsubscript{2} presentan compensaciones y cotas activas;
    \item si el observador comienza después del agotamiento efectivo del O\textsubscript{2}, puede ser razonable
    estudiar la aproximación $f_{\mathrm{ferm}}\approx1$.
\end{enumerate}

\subsection{Decisión pendiente sobre $S$}

Antes de implementar el EKF de cinco estados hay que definir una de estas opciones:
\begin{enumerate}
    \item mantener $G$ y $F$ como estados separados en una primera versión del observador;
    \item usar una trayectoria nominal $\rho_G(t)$ para reconstruir $G$ y $F$ desde $S$;
    \item desarrollar y validar una cinética agregada de azúcar que cierre directamente $\dot S$.
\end{enumerate}

La opción 1 es la más fiel a la implementación actual; la opción 3 es la más limpia para el
vector final $[S,N,X,E,C_d]^\mathsf{T}$, pero requiere una reducción adicional que todavía no ha
sido validada.

\section{Limitaciones que deben acompañar el modelo}

\begin{itemize}
    \item $C_d(0)=0$ es un supuesto de la capa CO\textsubscript{2}, no una medición de CO\textsubscript{2} disuelto.
    \item El $O_2$ actual es un estado latente efectivo y no debe interpretarse todavía como una trayectoria fisiológica validada.
    \item No existe un estado explícito de \emph{headspace}; parte de la cola temporal puede corresponder a fenómenos no incluidos.
    \item $k_{\mathrm{release}}$, $s_{\mathrm{sat}}$, $g_m$ y los parámetros de activación pueden compensarse.
    \item Los parámetros CO\textsubscript{2} actuales son específicos de matriz; no están demostrados como universales MS/MN.
    \item La agregación $S=G+F$ no es una reducción cerrada sin una hipótesis adicional sobre la partición de azúcares.
    \item La respuesta gradual al pulso de N en la capa CO\textsubscript{2} utiliza una simulación contrafactual sin pulso; una formulación
    de EKF completamente autocontenida deberá representar esta memoria de manera explícita o justificar su simplificación.
\end{itemize}

\section{Fuentes de implementación}

Las ecuaciones y valores de este documento se reconstruyeron de las siguientes rutas del branch
\texttt{ctorrealba\_fermentation}:
\begin{itemize}
    \item \texttt{fermentation\_model/shared/run\_new\_must\_glycerol\_estimability\_doe.py}
    \item \texttt{fermentation\_model/shared/run\_secondary\_joint\_campaign\_doe.py}
    \item \texttt{fermentation\_model/laboratory\_2026/run\_co2\_matrix\_cross\_validation\_2026.py}
    \item \texttt{fermentation\_model/pilot\_2025/run\_pilot\_2025\_co2\_solubility\_integrated\_doe.py}
    \item \texttt{fermentation\_model/laboratory\_2026/results/estimability\_historical\_natural/theta.csv}
    \item \texttt{fermentation\_model/laboratory\_2026/results/estimability\_historical\_synthetic\_plus\_lot2/theta\_by\_case.csv}
    \item \texttt{fermentation\_model/laboratory\_2026/results/co2\_matrix\_cross\_validation\_2026/fit\_parameters.csv}
\end{itemize}

\paragraph{Nota.}
Este documento debe entenderse como una \textbf{formalización del estado actual del trabajo}, no como
una afirmación de que todos los parámetros sean propiedades físicas identificadas de manera única.

\end{document}
